% ---------------------------------------------------------------------
% Introducción
% ---------------------------------------------------------------------

\chapter{Introducción}
\label{cap:introduccion}

\begin{Resumen}
Este capítulo presenta el contexto del Trabajo de Fin de Grado, la motivación técnica del proyecto y el alcance inicial de la memoria. El trabajo se plantea como la construcción e integración de una plataforma aérea propia para capturar datos reales y estudiar la viabilidad de una navegación asistida por visión en escenarios en los que la señal GNSS no sea fiable o no esté disponible.
\end{Resumen}

\section{Contexto}
\label{sec:introduccion-contexto}

Los vehículos aéreos no tripulados han pasado de ser plataformas principalmente orientadas al pilotaje remoto a convertirse en sistemas robóticos capaces de ejecutar misiones semiautónomas. Esta evolución se apoya en controladores de vuelo cada vez más completos, sensores de bajo coste, enlaces de comunicación robustos y computadoras de borde que permiten procesar información a bordo.

En la mayoría de aplicaciones civiles, la navegación de un dron depende de sistemas GNSS para estimar su posición global. Esta dependencia es suficiente para tareas de inspección, cartografía o seguimiento de rutas en condiciones normales, pero representa una limitación cuando la señal se degrada, se bloquea o sufre interferencias. En estos casos, disponer de una fuente alternativa de localización puede aumentar la robustez del sistema y permitir operaciones más seguras.

Una de las alternativas más prometedoras es la navegación basada en visión. En lugar de depender únicamente de coordenadas satelitales, el dron puede utilizar imágenes capturadas durante el vuelo y compararlas con una base de datos georreferenciada. Esta idea se conoce como Geolocalización Visual Cruzada o \textit{Cross-View Geo-Localization} (CVGL), ya que compara imágenes tomadas desde puntos de vista diferentes, por ejemplo una cámara embarcada y un mosaico satelital.

\section{Motivación}
\label{sec:introduccion-motivacion}

La motivación principal del proyecto es construir una plataforma real que permita pasar de la teoría de la geolocalización visual a un banco de pruebas operativo. El interés del trabajo no se limita al algoritmo de visión, sino que incluye la integración completa del sistema: estructura del dron, propulsión, controlador de vuelo, cámara, computadora de borde, comunicaciones, captura de datos y flujo de evaluación.

El uso de un dron construido manualmente permite controlar las decisiones de diseño y comprender las limitaciones prácticas que aparecen en una plataforma ligera: vibraciones, autonomía, distribución de masas, alimentación eléctrica, latencia de captura, carga computacional y fiabilidad de los enlaces de comunicación. Estas limitaciones son especialmente importantes cuando se pretende ejecutar procesamiento visual a bordo.

\section{Problema planteado}
\label{sec:introduccion-problema}

El problema que aborda este trabajo puede formularse como la integración de un sistema aéreo capaz de capturar imágenes georreferenciadas durante una misión real y utilizar esas imágenes para estimar la posición del dron mediante comparación visual con un mapa de referencia. Para ello es necesario resolver tres subproblemas:

\begin{itemize}
\item diseñar y ensamblar una plataforma aérea estable y suficientemente capaz para transportar cámara, computadora de borde y comunicaciones;
\item configurar el controlador de vuelo y el enlace MAVLink para obtener telemetría sincronizada con las imágenes;
\item preparar un flujo de CVGL que permita comparar los fotogramas de vuelo con una galería satelital y medir el error de localización.
\end{itemize}

\section{Alcance del trabajo}
\label{sec:introduccion-alcance}

El alcance inicial del proyecto se centra en la integración y validación experimental de la plataforma. El dron se ha construido a partir de un chasis tipo F450, motores brushless, hélices de 10 pulgadas, batería LiPo 4S, controlador de vuelo H743 con ESC, módulo GPS con brújula, cámara USB de obturador global y una Raspberry Pi 3 B+ como computadora de borde (ver \autoref{fig:dron-general}).

En la versión actual, la Raspberry Pi se utiliza principalmente para capturar imágenes y guardar telemetría asociada. La ejecución completa de la geolocalización visual se plantea inicialmente fuera de bordo, en un ordenador con más capacidad de cálculo, debido a que el procesamiento con descriptores visuales y emparejamiento local resulta demasiado pesado para la Raspberry Pi disponible. Esta decisión permite validar primero la adquisición de datos y el algoritmo de localización antes de abordar una optimización para tiempo real embarcado.

\begin{figure}[htbp]
\centering
\includegraphics[width=0.75\textwidth]{dron_general_anotado.png}
\caption{Vista general de la plataforma aérea desarrollada, con los componentes principales señalados: GPS, Raspberry Pi, stack FC/ESC, BEC, batería y cámara.}
\label{fig:dron-general}
\end{figure}

\section{Metodología general}
\label{sec:introduccion-metodologia}

La metodología seguida en el proyecto se estructura en fases. En primer lugar se seleccionan los componentes y se monta la plataforma aérea. A continuación se configura el controlador de vuelo mediante herramientas de estación de tierra, se calibra la electrónica principal y se comprueba la estabilidad del sistema. Después se integra la Raspberry Pi con la cámara y con el flujo de telemetría MAVLink, de forma que cada sesión de vuelo genere un conjunto de imágenes y un archivo CSV con la información temporal y de navegación.

Una vez obtenidos los datos de vuelo, se descarga un mosaico satelital de la zona, se divide en parches georreferenciados y se construye una galería visual. Cada imagen capturada por el dron se compara con esta galería mediante una fase de recuperación global y una fase posterior de emparejamiento local. Finalmente se calcula la posición estimada y se compara con la telemetría GNSS registrada durante el vuelo, que se utiliza como referencia para cuantificar el error.

\section{Estructura de la memoria}
\label{sec:introduccion-estructura}

La memoria se organiza de la siguiente forma. El \autoref{cap:objetivos-competencias} (\nameref{cap:objetivos-competencias}) recoge los objetivos y competencias asociadas. El \autoref{cap:estado-arte} (\nameref{cap:estado-arte}) resume los conceptos necesarios sobre navegación de drones, MAVLink y geolocalización visual. El \autoref{cap:diseno-sistema} (\nameref{cap:diseno-sistema}) describe la arquitectura general del sistema. Los capítulos \ref{cap:integracion-hardware} (\nameref{cap:integracion-hardware}) y \ref{cap:integracion-software} (\nameref{cap:integracion-software}) detallan la integración física y lógica de la plataforma. El \autoref{cap:geolocalizacion-visual} (\nameref{cap:geolocalizacion-visual}) presenta el flujo de CVGL utilizado. El \autoref{cap:experimentos-resultados} (\nameref{cap:experimentos-resultados}) documenta los vuelos, métricas y resultados obtenidos. El \autoref{cap:planificacion-presupuesto} (\nameref{cap:planificacion-presupuesto}) incluye la planificación y el presupuesto, y por último el \autoref{cap:conclusiones} (\nameref{cap:conclusiones}) recoge las conclusiones y líneas futuras.

% ---------------------------------------------------------------------